
\documentclass[11pt,a4paper]{article}

% ---------- Кодировка и язык ----------
\usepackage[T2A]{fontenc}
\usepackage[utf8]{inputenc}
\usepackage[english,russian]{babel}

% ---------- Оформление ----------
\usepackage[left=2cm,right=2cm,top=1.8cm,bottom=1.8cm]{geometry}
\usepackage{microtype}
\usepackage{setspace}
\setstretch{1.08}
\usepackage{parskip}

% ---------- Математика и таблицы ----------
\usepackage{amsmath,amssymb}
\usepackage{booktabs}
\usepackage{array}
\usepackage{tabularx}

% ---------- Ссылки ----------
\usepackage[hidelinks]{hyperref}
\usepackage{xurl}

% ---------- Списки ----------
\usepackage{enumitem}
\setlist[itemize]{leftmargin=*,nosep,topsep=3pt}
\setlist[enumerate]{leftmargin=*,nosep,topsep=3pt}

% ---------- Заголовки ----------
\usepackage{titlesec}
\titleformat{\section}
  {\normalfont\bfseries\normalsize}
  {\thesection.}{0.5em}{}
\titlespacing*{\section}{0pt}{10pt}{4pt}

% ---------- Колонтитулы ----------
\usepackage{fancyhdr}
\pagestyle{fancy}
\fancyhf{}
\fancyfoot[C]{\thepage}
\renewcommand{\headrulewidth}{0pt}

% ---------- Пользовательские команды ----------
\newcommand{\todo}[1]{\textbf{[TODO: #1]}}
\newcommand{\acc}{\mathrm{Acc}}
\newcommand{\pp}{\text{ п.п.}}

\usepackage{float}

\begin{document}

\begin{center}
    {\Large\bfseries Воспроизводимость и оценка
    эффективности ConvNeXt}\\[4pt]
    {\normalsize Классификация изображений малого разрешения}\\[6pt]
    \textbf{Промежуточная точка проекта}\\[4pt]
    Анастасия Оленева\\
    НИУ ВШЭ
\end{center}

\vspace{-0.5em}

\section{Статья и основная гипотеза}

\textbf{Статья:}
Feng J. et al.
\textit{Enhancing ConvNeXt for Efficient Small-Size
Image Classification}, 2025.

\textbf{Код:}
\url{https://github.com/southerly7/Proposed}

\textbf{Гипотеза.}
Авторы утверждают, что адаптация архитектуры ConvNeXt к изображениям малого разрешения, в частности использование двухэтапного patch embedding (Twice-Patch Embedding), модифицированных свёрточных блоков и дополнительных механизмов нормализации и внимания, позволяет повысить точность классификации при относительно небольшом количестве параметров. 


\section{Ядро воспроизведения}

\textbf{Датасет:} CIFAR-100, изображения $32\times32$,
100 классов.

\textbf{Основная метрика:} Top-1 Accuracy (\%).

Для воспроизведения выбирается результат предложенной модели из таблицы 3 оригинальной статьи.

\begin{table}[H]
\centering

\begin{tabular}{l l l}
Модель & Параметры & Accuracy \\
ConvNeXt-T & 27,8 млн & 67,58\% \\
Conv2NeXt-B & 7,3 млн & 83,82\% \\
Предложенная модель & 7,3 млн & 85,16\% \\

\end{tabular}

\end{table}

 
Основной задачей является воспроизведение точности предложенной модели (85,16\%) с использованием авторского кода.

Дополнительно проверяется, насколько результат устойчив к случайной инициализации: планируется выполнить три независимых запуска с различными seed.

\textbf{Критерий успеха:}
средняя Top-1 Accuracy отличается от опубликованной не более чем на 1 процентный пункт. Поскольку авторы не приводят стандартное отклонение по seed, этот порог принимается предварительно и будет дополнительно оценён по фактическому разбросу результатов. 


\section{Запуск оригинального кода}

Первый запуск выполнен на основе оригинального
репозитория авторов без изменения архитектуры модели
и алгоритма обучения. Из-за несовместимости
с современными версиями библиотек потребовались
небольшие технические исправления.

\begin{itemize}
    \item Commit hash: 0c57846d62c25ac548d7810d425f00b9299a8e54 
    \item Python: 3.13.16
    \item \item PyTorch: 2.14.1+cu130
    \item torchvision: 0.29.1+cu130
    \item timm: 1.0.31
    \item CUDA: 13.0
    \item GPU: NVIDIA Tesla T4
    \item Время обучения: 5 мин 37 с
          ($\approx 0.094$ GPU-ч)
    \item Полученная accuracy: 11.91\%
          (после 2 эпох)
    \item Статус: проверочный запуск успешен,
          потребовались исправления совместимости.
\end{itemize}

Указанные авторами версии:
PyTorch 1.8.0, torchvision 0.9.0,
timm 0.3.2.

Для запуска были внесены следующие изменения:
\begin{itemize}
    \item Удалены импорты недоступных оптимизаторов
          из \texttt{timm} (в эксперименте
          используется AdamW из PyTorch).
    \item Устаревший импорт
          \texttt{torch.\_six.inf}
          заменён на \texttt{math.inf}.
    \item Установлена дополнительная зависимость
          \texttt{tensorboardX}.
    \item Добавлена обработка аргументов
          \texttt{pretrained\_cfg} и
          \texttt{cache\_dir} при создании модели.
\end{itemize}

Проверочный запуск модели \texttt{proposed\_base}
на CIFAR-100 выполнен с seed 42 в течение
двух эпох. Количество параметров модели
составило 7.41 млн. Top-1 Accuracy увеличилась
с 7.28\% до 11.91\%, а train loss уменьшился
с 4.5738 до 4.3965. Сохранение checkpoint
проверено успешно.

Полное обучение на 300 эпох пока не проводилось.
Его предварительная вычислительная стоимость
оценивается примерно в 14 GPU-часов.
Результаты проверочного запуска не сравниваются
с заявленной авторами accuracy 85.16\%,
поскольку использовалось сокращённое обучение
без warmup.

\section{Проверяемые недостатки}

\textbf{1. Неизолированный вклад модифицированного stem }

В статье предлагается адаптированный механизм Twice-Patch Embedding, позволяющий лучше сохранять пространственную информацию при обработке изображений малого разрешения. Однако изменение stem рассматривается вместе с другими модификациями архитектуры, что затрудняет оценку его независимого вклада в итоговую точность. 

\textit{Проверка:}
провести абляционный эксперимент, в котором остальные компоненты сети остаются неизменными, а изменяется только входной блок. Планируется сравнить стандартный stem ConvNeXt с ядром \textbackslash{}(4\textbackslash{}times4\textbackslash{}) и шагом 4, адаптированный одноэтапный stem и предложенный авторами Twice-Patch Embedding. 

Для каждого варианта будут измерены Top-1 Accuracy, число параметров и FLOPs. Поскольку различные stem могут изменять разрешение промежуточных карт признаков, эти различия будут учитываться при интерпретации результатов. 

\textit{Ожидание:}
адаптированный stem обеспечит более высокую точность на CIFAR-100 за счёт сохранения пространственной информации на ранних стадиях. Предполагаемый прирост составляет 1–3 процентных пункта относительно стандартного варианта. 

\medskip
\textbf{2. Недостаточно строгое сравнение с ResNet }

В статье предложенная архитектура сравнивается с несколькими моделями, однако число параметров и условия обучения различаются. Поэтому нельзя однозначно определить, связано ли преимущество с архитектурными решениями или с различиями в вычислительном бюджете и процедуре обучения. 

\textit{Проверка:}
обучить ResNet-18 с подобранной шириной каналов таким образом, чтобы число параметров отличалось от предложенной ConvNeXt не более чем на 5\%.

Для обеих моделей будут использоваться одинаковые данные, аугментации, оптимизатор, scheduler, количество эпох и протокол оценки. Дополнительно сравниваются FLOPs и время обучения.

 

\textit{Ожидание:}
улучшенная ConvNeXt сохранит преимущество над ResNet при сопоставимом количестве параметров, однако разница в точности окажется меньше, чем при сравнении моделей в исходной статье. 

\medskip
\textbf{3. Потенциальная утечка данных при выборе checkpoint }

В оригинальной реализации тестовая выборка CIFAR-100 используется для оценки качества после каждой эпохи и выбора лучшего checkpoint. Такой подход может приводить к оптимистичному смещению итоговой метрики, поскольку тестовые данные косвенно участвуют в выборе модели. 

\textit{Проверка:}
выделить из обучающей выборки независимый validation set (45 000 изображений для обучения и 5 000 для валидации). Лучший checkpoint будет выбираться исключительно по validation accuracy, а test set использоваться только для итоговой оценки.

Будут сопоставлены результаты оригинальной и исправленной процедур. При анализе будет учитываться влияние уменьшения обучающей выборки.

 

\textit{Ожидание:}
после исправления протокола итоговая accuracy может несколько снизиться. Предварительно ожидается разница порядка 0,2–1,0 процентного пункта. 

\section{Улучшение и расширение}

В качестве дополнительного эксперимента предлагается исследовать, сохраняется ли преимущество адаптированного stem при уменьшении размера модели.

Для этого будет использована уменьшенная конфигурация ConvNeXt-Tiny (около 3,3 млн параметров), которая сравнивается с основной конфигурацией Base (около 7,3 млн параметров).

Для каждой конфигурации будут рассмотрены варианты со стандартным и адаптированным stem.

Основная исследуемая величина:

$\Delta_{stem} = Accuracy_{adaptive} - Accuracy_{standart}$


\textbf{Ожидание:}
уменьшенная модель продемонстрирует более низкую абсолютную точность, однако преимущество адаптированного stem сохранится. Величина прироста может зависеть от числа параметров и ёмкости модели.

Этот эксперимент является дополнительным и будет выполняться при наличии вычислительных ресурсов.


\section{Бюджет и план}

Время полного обучения будет оценено
после профилирования первых 5--10 эпох.

Пусть $T_B$ --- время одной эпохи Base,
$T_T$ --- время одной эпохи Tiny.

Тогда оценка времени полного запуска:
\[
C_B \approx 300T_B,
\qquad
C_T \approx 300T_T.
\]

\begin{table}[h]
\centering
\small
\caption{Предварительный план экспериментов}
\begin{tabularx}{\linewidth}{Xcc}
\toprule
Эксперимент & Запуски & Эпохи \\
\midrule
Оригинальный код & 1 & 300 \\
Base: полный и без модулей, 3 seed & 6 & 300 \\
Проверка корректной оценки & 1--3 & 300 \\
Tiny: scaling-эксперимент & 2 & 300 \\
GRN-only (опционально) & 1 & 300 \\
\bottomrule
\end{tabularx}
\end{table}

Оценка GPU-часов:
по результатам проверочного запуска
на NVIDIA Tesla T4 среднее время одной эпохи
составило около 2.8 мин.
Таким образом, полное обучение модели Base
в течение 300 эпох потребует приблизительно
14 GPU-часов без учёта накладных расходов..

Проект выполняется индивидуально.

При недостатке ресурсов сначала исключается
GRN-only, затем scaling-эксперимент.
Приоритет сохраняется за воспроизведением
основной модели и аудитом оценки.

\section{Декларация использования LLM}

При подготовке плана использовался ChatGPT
(OpenAI) для разбора статьи, обсуждения
экспериментального дизайна, анализа
исходного кода и редактирования текста.

Проверка кода, запуск экспериментов,
анализ полученных данных и формулирование
итоговых выводов выполняются автором.
Все численные результаты получены
непосредственно в экспериментах.

\end{document}

